\section{新前沿 I：生成式 AI 与 Diffusion Models}

\subsection{生成建模的发展脉络}

在 Lecture 4 中，课程介绍了两种经典的生成模型——VAE（Variational Autoencoder）和 GAN（Generative Adversarial Network）。这两类模型虽然开创性地实现了从数据中学习生成新样本的能力，但它们面临着几个关键的局限：

\begin{figure}[H]
\centering
\includegraphics[width=0.85\textwidth]{slides-images/slide-48.jpg}
\caption{VAE/GAN 的局限与挑战\protect\footnotemark}
\end{figure}
\footnotetext{Slides 第48页。}

\begin{knowledgebox}{VAE 和 GAN 的主要问题}
\textbf{局限性（Limitations）}：
\begin{itemize}
    \item \textbf{Mode collapse}：模型倾向于只生成``平均''样本，丧失多样性
    \item \textbf{分布外生成困难}：难以生成训练分布之外的新颖样本
    \item \textbf{训练困难}：特别是 GAN 的对抗训练过程极不稳定
\end{itemize}
\textbf{挑战（Challenges）}：稳定性、效率、生成质量、新颖性
\end{knowledgebox}

\subsection{Diffusion Models 的核心思想}

Diffusion Models 提出了一种与 VAE/GAN 截然不同的生成范式：

\begin{figure}[H]
\centering
\includegraphics[width=0.85\textwidth]{slides-images/slide-50.jpg}
\caption{VAE/GAN 的 ``one-shot'' 生成 vs. Diffusion 的迭代去噪生成\protect\footnotemark}
\end{figure}
\footnotetext{Slides 第50页。}

\begin{importantbox}{Diffusion Models 的关键创新}
\begin{itemize}
    \item \textbf{VAE/GAN}：从低维隐变量\textbf{一步到位}（one-shot）地生成完整样本
    \item \textbf{Diffusion Models}：通过\textbf{迭代地}（iteratively）逐步去噪来生成样本
\end{itemize}
这种迭代策略将困难的生成任务分解为一系列简单的去噪步骤，从而显著提升了生成质量。
\end{importantbox}

\subsection{前向加噪过程 (Forward Noising)}

Diffusion 模型的第一步是构建训练数据——通过一个前向加噪过程（Forward Noising Process）。

\begin{figure}[H]
\centering
\includegraphics[width=0.85\textwidth]{slides-images/slide-51.jpg}
\caption{扩散过程：前向加噪与反向去噪\protect\footnotemark}
\end{figure}
\footnotetext{Slides 第51页。引用 Sohl-Dickstein+ ICML 2015; Ho+ NeurIPS 2020。}

\begin{figure}[H]
\centering
\includegraphics[width=0.85\textwidth]{slides-images/slide-53.jpg}
\caption{前向加噪过程：逐步添加噪声\protect\footnotemark}
\end{figure}
\footnotetext{Slides 第53页。}

前向过程的关键步骤：
\begin{enumerate}
    \item 给定一张训练图像 $x_0$
    \item 采样一个随机噪声模式
    \item 按预定的时间表（schedule）逐步增加噪声比例：
    \begin{itemize}
        \item $T=0$：100\% 图像，0\% 噪声
        \item $T=1$：75\% 图像，25\% 噪声
        \item $T=2$：50\% 图像，50\% 噪声
        \item $T=3$：25\% 图像，75\% 噪声
        \item $T=4$：0\% 图像，100\% 噪声（纯随机噪声）
    \end{itemize}
\end{enumerate}

\begin{knowledgebox}{前向过程不需要训练}
前向加噪过程完全是确定性的（给定噪声 schedule），不涉及任何可学习参数。它的唯一目的是为反向去噪过程生成训练对——每一对 $(x_t, x_{t-1})$ 就是一个训练样本。
\end{knowledgebox}

\subsection{反向去噪过程 (Reverse Denoising)}

训练的核心任务是让神经网络学会反向过程——给定时刻 $T$ 的带噪图像，估计时刻 $T-1$ 的更干净的图像。

\begin{figure}[H]
\centering
\includegraphics[width=0.85\textwidth]{slides-images/slide-54.jpg}
\caption{反向去噪：学习从 $T$ 估计 $T-1$\protect\footnotemark}
\end{figure}
\footnotetext{Slides 第54页。引用 Ho+ NeurIPS 2020。}

\begin{importantbox}{Diffusion 训练目标}
给定时刻 $T$ 的带噪图像 $x_T$，训练一个神经网络 $f_\theta$ 来预测时刻 $T-1$ 的图像 $x_{T-1}$：
$$
\hat{x}_{T-1} = f_\theta(x_T, T)
$$
训练损失为预测图像与真实图像之间的差异（通常是 MSE 或预测噪声的 MSE）。这个网络在所有时间步上共享参数，时间步 $T$ 作为条件输入。
\end{importantbox}

\subsection{采样生成新图像}

训练完成后，生成新图像的过程就是从纯噪声开始，反复应用去噪网络：

\begin{figure}[H]
\centering
\includegraphics[width=0.85\textwidth]{slides-images/slide-61.jpg}
\caption{采样过程：从纯噪声逐步生成清晰图像\protect\footnotemark}
\end{figure}
\footnotetext{Slides 第61页。}

采样过程：$x_T \xrightarrow{f_\theta} x_{T-1} \xrightarrow{f_\theta} x_{T-2} \xrightarrow{f_\theta} \cdots \xrightarrow{f_\theta} x_0$

每一步去噪都只需要移除``一小部分''噪声，这比一步到位生成完整图像要简单得多——这正是 Diffusion Models 生成质量优越的根本原因。

\subsection{从图像到自然语言：Text-to-Image}

Diffusion Models 的一个重要应用是与自然语言处理的结合——文本生成图像（Text-to-Image Generation）。

\begin{figure}[H]
\centering
\includegraphics[width=0.85\textwidth]{slides-images/slide-64.jpg}
\caption{从自然语言描述生成图像\protect\footnotemark}
\end{figure}
\footnotetext{Slides 第64页。引用 Ramesh+ arXiv 2022。}

\begin{figure}[H]
\centering
\includegraphics[width=0.85\textwidth]{slides-images/slide-65.jpg}
\caption{Text-to-Image 生成的多样化示例\protect\footnotemark}
\end{figure}
\footnotetext{Slides 第65页。引用 OpenAI, Ramesh+ arXiv 2022。}

当今主流的 Text-to-Image 系统（如 DALL-E、Stable Diffusion、Midjourney）的图像生成骨干网络大多采用 Diffusion Model。用户输入的文本通过语言编码器转化为条件向量，引导去噪过程朝着符合文本描述的方向生成图像。

\subsection{超越图像：分子设计}

\begin{figure}[H]
\centering
\includegraphics[width=0.85\textwidth]{slides-images/slide-66.jpg}
\caption{Diffusion Models 在分子设计和蛋白质生成中的应用\protect\footnotemark}
\end{figure}
\footnotetext{Slides 第66页。}

Diffusion Models 的应用远不止于图像生成。在化学和生物学领域，同样的``从噪声到结构''的范式被用于：
\begin{itemize}
    \item \textbf{分子生成}：从 3D 噪声中逐步生成药物分子结构
    \item \textbf{蛋白质设计}：生成具有特定功能的新型蛋白质结构
\end{itemize}

\begin{knowledgebox}{Diffusion 的通用性}
Diffusion Models 的``加噪 $\rightarrow$ 去噪''框架具有高度通用性：只要数据可以被逐步``破坏''并``恢复''，这个框架就可以适用。这使得它成为当前最活跃的生成建模范式之一，跨越图像、视频、音频、3D 结构、分子等多个领域。
\end{knowledgebox}

\subsection{本章小结}

Diffusion Models 通过将生成任务分解为多步迭代去噪，克服了 VAE/GAN 在生成质量、稳定性和多样性方面的局限。其核心思想包括：前向加噪构建训练数据，反向去噪学习生成过程，最终从纯噪声采样出高质量样本。该范式已广泛应用于 Text-to-Image、视频生成、分子设计等领域。

%% =====================================================
